\subsection*{3.4 Transforming to a monic polynomial}
Given any polynomial $p(\mathbf z)$ in variables $\mathbf z=(z_1,\ldots,z_n)$, there is a standard trick to make it monic in a new variable $x$ by applying a linear transformation on the variables: for $\boldsymbol\alpha=(\alpha_1,\ldots,\alpha_n)\in\mathbb F^n$, let
\[\tau_{\boldsymbol\alpha}:z_i\mapsto\alpha_ix+z_i,\]
for $i=1,\ldots,n$. Let $p_{\boldsymbol\alpha}(x,\mathbf z)$ be the resulting polynomial. Note that $p$ and $p_{\boldsymbol\alpha}$ have the same degree. We show that $p_{\boldsymbol\alpha}(x,\mathbf z)$ is monic in $x$, for a random transformation $\tau_{\boldsymbol\alpha}$.
\begin{ntlemma}[Author Lemma 7: Transformation to monic]
Let $p(\mathbf z)$ be polynomial of total degree $d$. Let $S\subseteq\mathbb F$ be a finite set. For $\boldsymbol\alpha\in S^n$ picked independently and uniformly at random,
\[\Pr[p_{\boldsymbol\alpha}(x,\mathbf z)\text{ is monic in }x]\geq1-\frac d{|S|}.\]
\end{ntlemma}
\begin{proof}
Consider the terms of degree $d$ in $p$. Let $\boldsymbol\beta=(\beta_1,\ldots,\beta_n)$ such that $|\boldsymbol\beta|=\sum_{i=1}^n\beta_i=d$. We denote the term $\mathbf z^{\boldsymbol\beta}=z_1^{\beta_1}\cdots z_n^{\beta_n}$. Then the homogeneous component of degree $d$ in $p$ can be written as $a_d(\mathbf z)=\sum_{|\boldsymbol\beta|=d}c_{\boldsymbol\beta}\mathbf z^{\boldsymbol\beta}$. Note that $a_d$ is a nonzero polynomial.
Now consider the transformed polynomial $p_{\boldsymbol\alpha}$. We have $\deg_x(p_{\boldsymbol\alpha})=d$ and the coefficient of $x^d$ in $p_{\boldsymbol\alpha}$ is $a_d(\boldsymbol\alpha)=\sum_{|\boldsymbol\beta|=d}c_{\boldsymbol\beta}\boldsymbol\alpha^{\boldsymbol\beta}$. When we pick $\boldsymbol\alpha$ at random,$a_d(\boldsymbol\alpha)$ will be a nonzero constant with probability $\geq1-d/|S|$ by the PIT-Theorem, and in this case $p_{\boldsymbol\alpha}(x,\mathbf z)$ is monic in $x$.
\end{proof}
Given an ABP of size $s$ that computes $p(\mathbf z)$, we can construct another ABP of size $3s$ that computes $p_{\boldsymbol\alpha}(x,\mathbf z)$. For the new ABP replace edge labeled by $z_i$ by the ABP computing $\alpha_ix+z_i$. For each old edge, this requires adding two new edges with labels $\alpha_i$ and $x$.
\subsection*{3.5 Handling the starting point of Hensel lifting}
After doing the above pre-processing steps on the given polynomial $p(\mathbf z)$, we call the transformed polynomial $f(x,\mathbf z)$. We can assume that $f$ of degree $d$ can be factorized as $f=gh$, where $g$ and $h$ are coprime and $g$ is irreducible. In the first step of Hensel lifting, we factorize the univariate polynomial $f(x,0,\ldots,0)\equiv f(x,\mathbf z)\pmod{\mathbf z}$. Now, clearly we have the factorization $f(x,0,\ldots,0)=g(x,0,\ldots,0)h(x,0,\ldots,0)$, but these two factors might not be coprime. In this case we do another transformation.
\paragraph{Author Remark}
Although it would suffice for our purpose to start with two coprime factors, the transformation below produces one irreducible factor.
Let $g_0$ be an irreducible factor of $g(x,0,\ldots,0)$. Then we have for some univariate polynomial $h'_0(x)$ and for $h_0(x)=h'_0(x)h(x,0,\ldots,0)$,
\[g\equiv g_0h'_0\pmod{\mathbf z},\qquad f\equiv g_0h_0\pmod{\mathbf z}.\]
We want that $g_0$ is coprime to $h'_0$ and $h_0$. Directly, this might not be the case because all factors of $f(x,0,\ldots,0)$ might have multiplicity $>1$. However, we argue how to ensure this after a random shift $\boldsymbol\alpha$ of $f$. That is, we consider the function $f(x,\mathbf z+\boldsymbol\alpha)$
\begin{enumerate}
\item First, we show how to achieve that $g_0$ is coprime to $h'_0$.
Since $g$ is irreducible, it is also square-free, and hence,$\gcd(g,\partial g/\partial x)=1$. By Lemma 2, the resultant $r(\mathbf z)=\Res_x(g,\partial g/\partial x)$ is a polynomial of degree $\leq4d^2$ and $r(\mathbf z)\ne0$. Hence, at a random point $\boldsymbol\alpha\in[8d^2]^n$, we have $r(\boldsymbol\alpha)\ne0$ with high probability. At such a point $\boldsymbol\alpha$, we have that $g(x,\boldsymbol\alpha)$ is square-free. Therefore,$g(x,\mathbf z)$ is square-free modulo $(\mathbf z-\boldsymbol\alpha)$, or, equivalently,$g(x,\mathbf z+\boldsymbol\alpha)$ is square-free modulo $\mathbf z$. Hence, when we define $g_0$ and $h'_0$ from $g(x,\mathbf z+\boldsymbol\alpha)$ instead of $g(x,\mathbf z)$, they will be coprime.
\item Similarly, we can achieve that $g_0$ is coprime to $h_0$. By the first item, it now suffices to get $g_0$ coprime to $h(x,0,\ldots,0)$.
For showing this, we use that $g_0$ is coprime to $h'_0$ and prove that $g(x,0,\ldots,0)$ is coprime to $h(x,0,\ldots,0)$. Consider the resultant of $g$ and $h$ w.r.t.$x$, the polynomial $r'(\mathbf z)=\Res_x(g,h)$ has degree $\leq4d^2$. Since $g$ and $h$ are coprime,$r'(\mathbf z)\ne0$. Hence, at a random point $\boldsymbol\alpha\in[8d^2]^n$, we have $r'(\boldsymbol\alpha)\ne0$ with high probability, and hence $g(x,\boldsymbol\alpha)$ and $h(x,\boldsymbol\alpha)$ are coprime univariate polynomials. Therefore,$g(x,\mathbf z)$ and $h(x,\mathbf z)$ are coprime modulo $(\mathbf z-\boldsymbol\alpha)$, or, equivalently,$g(x,\mathbf z+\boldsymbol\alpha)$ and $h(x,\mathbf z+\boldsymbol\alpha)$ are coprime modulo $\mathbf z$.
\end{enumerate}
Combining the two items, a random point $\boldsymbol\alpha\in[8d^2]^n$ will fulfill both properties with high probability. So instead of factoring $f(x,\mathbf z)$, we do a coordinate transformation $\mathbf z\mapsto\mathbf z+\boldsymbol\alpha$ and factor $f(x,\mathbf z+\boldsymbol\alpha)$ instead. From these factors, we easily get the factors of $f(x,\mathbf z)$ by inverting the transformation. Note also that when $f(x,\mathbf z)$ is monic in $x$, the same holds for $f(x,\mathbf z+\boldsymbol\alpha)$.
In the next section, we do another transformation on the input polynomial. We apply a map on the variables that maps $x$ to $x$ and $z_i$ is mapped to $yz_i$, for a new variable $y$ and $i=1,\ldots,n$. Then we factorize the transformed polynomial modulo $y$. Note that in this case, going modulo $y$ has the same effect of going modulo $\mathbf z$. So we can use the above argument to ensure the starting condition for Hensel lifting is satisfied.
\subsection*{3.6 Reducing multivariate factoring to the bivariate case}
Factoring multivariate polynomials can be reduced to the case of bivariate polynomials (see~\cite{ref12}). Let $x,y$ and $\mathbf z=(z_1,\ldots,z_n)$ be variables and let $f(x,\mathbf z)$ be the given polynomial. With $f\in\mathbb F[x,\mathbf z]$, we associate the polynomial $\widehat f\in\mathbb F[x,y,\mathbf z]$ defined by
\[\widehat f(x,y,\mathbf z)=f(x,yz_1,\ldots,yz_n).\]
The point now is to consider $\widehat f$ as a polynomial in $\mathbb F[\mathbf z][x,y]$, that is, as a bivariate polynomial in $x$ and $y$ with coefficients in $\mathbb F[\mathbf z]$. We list some properties.
\begin{enumerate}\item $f(x,\mathbf z)=\widehat f(x,1,\mathbf z)$,\item $\deg(\widehat f)\leq2\deg(f)$,\item $f\text{ monic in }x\Longrightarrow\widehat f\text{ monic in }x$,\item $f=gh\Longrightarrow\widehat f=\widehat g\widehat h$,\item $\widehat f=g'h'\Longrightarrow f=g'(x,1,\mathbf z)h'(x,1,\mathbf z)$.\end{enumerate}
By property 4, factors of $f$ yield factors of $\widehat f$. The following lemma shows that also the irreducibility of the factors is maintained.
\begin{ntlemma}[Author Lemma 8]
Let $f$ be monic in $x$ and $g$ be a non-trivial irreducible factor of $f$. Then $\widehat g$ is a non-trivial irreducible factor of $\widehat f$.
\end{ntlemma}
\begin{proof}
By property 4 above,$\widehat g$ is a factor of $\widehat f$. We argue that $\widehat g$ is irreducible.
Let $\widehat g=uv$ be a factorization of $\widehat g$. By item 5 above, this yields a factorization of $g$ as $g=u(x,1,\mathbf z)v(x,1,\mathbf z)$. Since $g$ is monic in $x$, the same holds for $\widehat g$, and therefore also for factors $u,v$. Hence, either $u$ or $v$ must be constant, because otherwise they would provide a non-trivial factorization of $g$.
\end{proof}
Thus, to get an ABP for an irreducible factor $g$ of $f$, first we show that there is an ABP for the irreducible factor $\widehat g$. This yields an ABP for $g$ by substituting $g=\widehat g(x,1,\mathbf z)$.
Given an ABP $B_f$ of size $s$ for $f$, we get an ABP $B_{\widehat f}$ for $\widehat f$ by putting an edge labeled $y$ in series with every edge labeled $z_i$ in $B_f$, so that $B_{\widehat f}$ computes $yz_i$ at every place where $B_f$ uses $z_i$. Hence, the size of $B_{\widehat f}$ is at most $2s$.
\subsection*{3.7 Solving a linear system with polynomials as matrix entries}
We show how to solve a linear system $M\mathbf v=0$ for a polynomial matrix $M$ with entries from $\mathbb F[\mathbf z]$ given as ABPs. We are seeking for a nonzero vector $\mathbf v$. Note that such a $\mathbf v$ exists over the ring $\mathbb F[\mathbf z]$ iff it exists over the field $\mathbb F(\mathbf z)$.
Except for minor modifications, this follows from classical linear algebra. Kopparty, Saraf, and Shpilka~\cite[Lemma 2.6]{ref12} have shown the same result for circuits. The proof works as well for ABPs.
\begin{ntlemma}[Author Lemma 9: Solving linear systems, {\cite{ref12}}]
Let $M=(m_{i,j}(\mathbf z))_{i,j}$ be a polynomial matrix of dimension $k\times m$ and variables $\mathbf z=(z_1,\ldots,z_n)$, where the entries are polynomials $m_{i,j}\in\mathbb F[\mathbf z]$ that can be computed by ABPs of size $s$.
Then there is an ABP of size $\operatorname{poly}(k,m,s)$ computing a nonzero vector $\mathbf v\in\mathbb F[\mathbf z]^m$ such that $M\mathbf v=0$ (if it exists).
\end{ntlemma}
\begin{proof}
After swapping rows of $M$, we ensure that the $j\times j$ submatrix $M_j$ that consists of the first $j$ rows and the first $j$ columns has full rank, iteratively for $j=1,2,\ldots$.
For $j=1$ this means to find a nonzero entry in the first column and swap that row with the first row. If the first column is a zero-column, then $\mathbf v=(1\ 0\ \cdots\ 0)^T$ is a solution and we are done. To extend from $j$ to $j+1$, suppose we have ensured that $M_j$ has full rank. Now we search for a row from row $j+1$ on, such that after a swap with row $j+1$, the submatrix $M_{j+1}$ has full rank. This can be tested by Lemma 3. If no such row exists, then the process stops at $j$. If $j=m$ then $M$ has full rank and the zero vector is the only solution. Otherwise, assume the above process stops with $j<m$.
Now Cramer's rule can be used to find the unique solution $\mathbf u=(u_1\ u_2\ \cdots\ u_j)^T$ of the system
\[M_j\mathbf u=(m_{1,j+1}\ m_{2,j+1}\ \cdots\ m_{j,j+1})^T.\]
We have $u_i=\det M_j^i/\det M_j$, where $M_j^i$ is the matrix obtained by replacing the $i$-th column of $M_j$ by the vector $(m_{1,j+1}\ \cdots\ m_{j,j+1})^T$. Now, define
\[\mathbf v=(\det M_j^1\ \det M_j^2\ \cdots\ \det M_j^j\ -\det M_j\ 0\ \cdots\ 0)^T.\]
Then $\mathbf v$ is a solution to the original system. Its entries are determinants of matrices with entries computed by ABPs of size $s$. Hence, all the entries of $\mathbf v$ have ABPs of size $\operatorname{poly}(k,m,s)$.
\end{proof}
